\section{Adjunctions, finite products, and submodules}
\label{sec:furtherfoundations}
\begin{proposition}[The unit-fiber adjunction]\label{prop:unitfiberadjunction}
For a unital inverse semiring $T$, let $U_1(T)$ be its additive group fiber
containing $1$, with group zero $e_1=1+(-1)$. This is a unital ring, and
\[
 \operatorname{Hom}_{\mathrm{InvSemi}}(R_0,T)
 \cong\operatorname{Hom}_{\mathrm{Ring}}(R,U_1(T))
\]
naturally. Thus zero adjunction is fully faithful and left adjoint to
$U_1$. This is the composite of the zero-adjunction/zeroless adjunction
and the ring-of-scalars adjunction stated in \cite[Definition 3.6,
Section 3.1, Proposition 4.7 and Remark 4.14]{FGJM}; it differs from their
ordinary ring reflection.
\end{proposition}
\begin{proof}
The idempotent map is multiplicative and $e_1^2=e_1$, so the fiber is
closed under multiplication and contains $1$. Its group zero absorbs
multiplication within the fiber, since $e_1a=a+(-a)=e_1$. Thus it is a
ring. A semiring map preserves unique additive inverses, sends the old
zero of $R_0$ to $e_1$, and restricts to a ring map into this fiber. Conversely
such a ring map extends uniquely by sending the new zero to the global
zero of $T$. The constructions are inverse and natural. The adjunction
unit $R\to U_1(R_0)$ is an isomorphism, proving full faithfulness.
\end{proof}

\begin{corollary}\label{cor:no-sections}
There is no unital semiring map from any ring to $R_0$ or to $\B$.
In particular, the ring reflection $R_0\to R$ has no semiring section.
If $R\ne0$, there is no unital semiring map $\B\to R_0$ either.
\end{corollary}
\begin{proof}
The images of $1$ and $-1$ under a ring-to-semiring map would sum to the
new zero. This is impossible in either target, where a sum is the new
zero only when both summands are new zero, whereas $1$ is not new zero.
A map $\B\to R_0$ would give $1_R+1_R=1_R$, hence $1_R=0_R$.
\end{proof}

There is nevertheless a continuous map
$\Spec\Z_0\to\Spec\Z$ which preserves every ordinary prime and sends the
new generic point to $(0)$. For $n\ne0$, the inverse image of
$D_{\Z}(n)$ is $D_{\Z_0}(n)$; for $n=0$ it is empty. These basic opens
prove continuity. This topological collapse is not induced by a unital
map $\Z\to\Z_0$, since Corollary~\ref{cor:no-sections} excludes such a map.

\begin{proposition}\label{prop:finite-product-reflections}
For a finite set $I$, the universal ring reflection of $(R_0)^I$ is $R^I$,
with zero fiber the Boolean cube $\B^I$. Localization of a finite semiring
product at one tuple is the product of its coordinate localizations.
For a field $K$, this turns nonzero-amplitude coordinates into unchanged
$K_0$ factors, old-zero coordinates into $\B$, and new-zero coordinates
into terminal factors. Identifying an old-zero mask with a tuple of
nonzero amplitudes on the same selected coordinates Booleanizes exactly
those coordinates.
\end{proposition}
\begin{proof}
Every additive idempotent is killed by a ring-valued map. Adding the full
old-zero mask to a tuple replaces its absent coordinates by old zero,
without changing that map's value. Thus the map factors uniquely through
its tuple of amplitudes. The zero fiber is exactly the cube of old/new
zero choices, with union and intersection operations.

Coordinate fractions admit a common denominator exponent, the maximum of
finitely many exponents; equality witnesses likewise admit a common
maximum. This proves the finite-product localization assertion. Inverting
an old field unit has no effect, inverting the old zero identifies all old
scalars with $1$, and inverting the new zero gives the terminal semiring.
For the mask congruence, multiply the proposed relation by each coordinate
projector. At a selected coordinate it says $e=a$ with $a\ne0$, hence makes
$e$ invertible and collapses the old field to its Boolean value. The
coordinate decomposition proves universality and preserves every
unselected coordinate.
\end{proof}

\begin{proposition}\label{prop:diagram-consequences}
The additive-idempotent functor $M\mapsto eM$ is both left and right
adjoint to restriction of $\B$-semimodules along $R_0\to\B$.
The least group fiber $M_{\bot}$ is right adjoint to restriction of
$R$-modules along $R_0\to R$. The cancellative $R_0$-semimodules are
exactly these ordinary $R$-modules.

A subsemimodule is specified by a zero-preserving join-subsemilattice
$L'\subseteq L_M$ and submodules $N_l\subseteq M_l$ for $l\in L'$,
compatible with every transition between its labels. If $m_1,\ldots,m_n$
generate $M$, its support lattice is generated by $l_i=em_i$ and has at most
$2^n$ elements; its fiber over $l$ is generated by the transported $m_i$
with $l_i\le l$.
\end{proposition}
\begin{proof}
A map to an idempotent module satisfies $f(m)=f(em)$, proving the left
adjunction. A map from an idempotent module lands in the additive
idempotents, proving the right one. A map from an ordinary module sends
inverse pairs to pairs summing to the global zero, so lands in $M_{\bot}$;
conversely a linear map to that fiber is an $R_0$-map. In a cancellative
module the identity $m+em=m$ gives $em=0$, so its support has one point.
Its group fiber is the required ordinary module, and the converse is
immediate.

A subsemimodule contains its supports, is closed under $-1$, and is closed
under adding a larger support. It therefore yields exactly the indicated
submodules and transitions. Conversely those data are closed under the
sum formula at a join and under scalar action. Every generated element
is a sum of old scalar multiples of the $m_i$, after omitting new-zero
terms. Its support is the join of their $l_i$. Within a fixed fiber,
transport the terms to that join to express it as the stated linear
combination. The empty sum supplies the global zero; no assertion that
the entire least group vanishes is used.
\end{proof}

\section{The multiplicative skeleton and its limits}
\label{sec:multiplicativeskeleton}
\begin{proposition}\label{prop:contracted-presentation}
Let $M$ be the multiplicative monoid of $R_0$. Its contracted monoid semiring
\[
 \mathbb N[M]/([\tau]=0)
\]
modulo the relations $[a]+[b]=[a+b]$ for $a,b\in R$ is $R_0$.
If one instead starts with the ordinary uncontracted monoid semiring and
imposes only those relations and $[x]+[\tau]=[x]$, the assertion fails.
\end{proposition}
\begin{proof}
After contraction, every nonempty sum of old generators reduces to one
old-ring generator. Evaluation distinguishes all such generators and the
empty sum, giving the isomorphism. Without contraction, evaluate each
$[x]$ in the inner copy of $R_0$ in $(R_0)_0$. All the displayed relations
hold, but $[\tau]$ maps to the inner zero and the empty sum to the distinct
outer zero. Thus those relations do not force preservation of the source
additive zero by the proposed inverse map.
\end{proof}

\begin{proposition}\label{prop:monoid-spectrum-ring}
If $R$ is a UFD with infinitely many irreducible classes, the spectrum of
the multiplicative monoid of $R_0$ has infinite dimension. If $R$ is a
domain, its ordinary contracted integral monoid ring is
\[
 \mathbb Z[M]/([\tau])\cong\mathbb Z\times\mathbb Z[R\setminus\{0\}].
\]
This ordinary ring must not be confused with the inverse-semiring monoid
algebra $\mathbb Z_0[M]$ of Proposition~\ref{prop:monoid}.
\end{proposition}
\begin{proof}
For any set $E$ of irreducible representatives, the set consisting of the
two zeros and all old nonzero elements divisible by some member of $E$
is a proper monoid ideal with multiplicatively closed complement.
Different sets give the same inclusion relation as their associated
ideals. Increasing finite sets of irreducibles give arbitrarily long
prime chains.

In the contracted integral monoid ring, let $z=[0_R]$. It is a central
idempotent and $z[a]=z$ for every old nonzero $a$. Its factor is therefore
$\mathbb Z$. The complementary factor has basis $[a]-z$ indexed by old
nonzero $a$; the domain hypothesis makes their products exactly
$([a]-z)([b]-z)=[ab]-z$. This identifies the other factor with the stated
ordinary monoid ring.
\end{proof}

For a domain, an old nonzero element has exactly its ordinary ring
factorizations, so its units and irreducibles are classical. The old
zero $e$ is reducible, since $e=e^2$, and is divisibility-prime: $e\mid xy$
means $xy\in\{e,\tau\}$, forcing a factor into that set. Thus prime in the
divisibility sense does not imply irreducible. The semiring prime spectrum
proved earlier is distinct from the much larger monoid spectrum.

\begin{proposition}[Ideal finiteness and integer arithmetic]
\label{prop:integer-arithmetic}
For a commutative unital ring $R$, extension $J\mapsto J\cup\{\tau\}$
preserves products of ideals and their finite generating sets. The
semiring $R_0$ is Noetherian, respectively has every ideal principal,
exactly when $R$ has that property.

In $\Z_0$, the cancellative multipliers are exactly the old nonzero
integers; the units are $\pm1$, and associates differ by sign. Define a
prime element to be a nonunit different from $\tau$ satisfying the usual
divisibility implication. The irreducibles are $\pm p$ for rational
primes $p$, whereas the prime elements are those elements together with
$e$. The old zero is reducible and has no factorization into irreducibles.
The Pythagorean triples are precisely the ordinary integer triples,
\[
 (a,\tau,\pm a),\qquad(\tau,b,\pm b)\quad(a,b\in\Z),
 \qquad(\tau,\tau,\tau).
\]
\end{proposition}
\begin{proof}
Finite sums of products of old elements are computed in $R$, and an old
zero term is available from every pair of ring ideals. Thus
$(I\cup\{\tau\})(J\cup\{\tau\})=IJ\cup\{\tau\}$. The ideal generated by
old elements $r_i$ is $(r_i)_R\cup\{\tau\}$; the new zero generates
$\{\tau\}$. The ideal poset of $R_0$ is the ideal poset of $R$ with one
new least point, so ascending chains stabilize in one exactly when they
do in the other. These facts prove both finiteness assertions.

Multiplication by an old nonzero integer distinguishes old inputs from
$\tau$, and is cancellative on the old inputs by integer cancellation.
Multiplication by $\tau$ is constant, while $e\cdot1=e\cdot2$.
A product equal to $1$ uses old integer factors, giving exactly the
units $\pm1$. Mutual divisibility for old nonzero integers is the ordinary
integer relation, while $e$ and $\tau$ each have a singleton associate
class, as their sets of multiples are $\{e,\tau\}$ and $\{\tau\}$.
Thus both usual definitions of associates give precisely sign changes.
Every factorization of an old nonzero integer likewise
uses old nonzero integers, so its irreducibility and primality are the
ordinary ones. The multiples of $e$ are exactly $\{e,\tau\}$; a product
lies in this pair only if one factor does, proving that $e$ is prime.
It is reducible because $e=e^2$, and a product of old irreducibles is
old nonzero.

If both entries on the left of $x^2+y^2=z^2$ are old, so is $z$, and the
equation is the ordinary integer one. If exactly one is $\tau$, the
remaining equation is $a^2=z^2$ in $\Z$, hence $z=\pm a$. If both are
$\tau$, then $z=\tau$. Every listed triple satisfies the equation.
\end{proof}

Although $R_0$ has no zero divisors with respect to $\tau$, for every
nonzero $R$ the equality $e\cdot e=e\cdot1$ prevents cancellation by all
multipliers other than $\tau$. When $R=0$, the semiring is $\B$, whose
only nonzero multiplier is $1$ and is cancellative.

If a finite nonnegative multiplicative ideal norm extends a nonconstant
ordinary ideal norm $N$, it must vanish on both $\{\tau\}$ and
$\{e,\tau\}$. Each of these ideals is unchanged by multiplication by
$J\cup\{\tau\}$ for a nonzero ordinary ideal $J$. Choosing $N(J)\ne1$
gives $w=wN(J)$ and hence $w=0$. Merely assuming that ordinary quotients
are finite does not make their cardinalities multiplicative. For
$R=\Z[2i]$ and $\mathfrak m=(2,2i)$, the basis $1,2i$ gives
$[R:\mathfrak m]=2$, while $\mathfrak m^2=(4,4i)$ has index $8$.
This supplies a concrete counterexample to that omitted hypothesis.

\section{Residue fibers and internal zero adjunction}
\label{sec:residueinternal}
\begin{proposition}\label{prop:residue-fibers}
For a surjective ring map $A\to k$ with kernel $\mathfrak m$, the induced
map $A_0\to k_0$ has old-element fibers $\widetilde a+\mathfrak m$ and
new-zero fiber $\{\tau\}$. For a finite tuple $a\in k_0^I$ of support $J$,
allow every absent coordinate of $a$ to remain absent or become the old
zero. The inverse image of this support star is the disjoint union,
indexed by $E\subseteq I\setminus J$, of
\[
 \prod_{i\in J}(\widetilde a_i+\mathfrak m)
 \times\prod_{i\in E}\mathfrak m
 \times\prod_{i\notin J\cup E}\{\tau\}.
\]
If $k$ is a field, reduction modulo $\mathfrak m^2$ makes each fixed
stratum an affine torsor under $(\mathfrak m/\mathfrak m^2)^{J\cup E}$.
Moreover, for every ideal $I\subseteq A$,
\[
 \varprojlim_n(A/I^n)_0\cong(\varprojlim_n A/I^n)_0.
\]
\end{proposition}
\begin{proof}
Old elements preserve occupied support, and the old quotient fiber is
precisely the stated coset. Partitioning the choices of activated absent
coordinates gives the product formula, with disjointness detected by
support. Quotienting each coset by $\mathfrak m^2$ gives its translation
torsor under the $k$-module $\mathfrak m/\mathfrak m^2$. This module is the
cotangent space; no extra scheme of cubes is implied by the subset
indexing. In the inverse limit, compatibility forces all coordinates of
a sequence to be old or all to be new zero. The old sequences are the
ordinary ring limit, and operations agree coordinatewise.
\end{proof}

In projective coordinates the matrix swapping two positions is invertible
and exchanges $[e:1]$ with $[1:e]$, and $[\tau:1]$ with $[1:\tau]$.
It therefore identifies the corresponding pointed residue-star sets.
For a nonzero $h\in\mathfrak m$, writing $[1:h]=[h^{-1}:1]$ requires
extension to a field containing $h^{-1}$. A nonmonomial fractional-linear
matrix, including a translation by a nonzero scalar, has no global
invertible lift over $K_0$, since an invertible Boolean matrix acts as an automorphism of the Boolean cube and must permute its atoms.
Cancelling a common factor from a pair of polynomials gives the ordinary
local extension wherever the reduced pair is unimodular; the uncancelled
pair $(e,e)$ at a common zero is not a projective point and differs from
$(\tau,\tau)$.

\begin{proposition}[Internal zero adjunction]\label{prop:internalzero}
For a commutative semiring object $A$ in a topos, the coproduct $A\amalg1$
carries zero adjunction, naturally in $A$. It has amplitude map to $A$ and
support map to $1\amalg1$, and it commutes with every geometric
inverse-image functor. For a sheaf $\mathcal A$ on a space,
\[
 (\mathcal A\amalg1)(U)
 \cong\coprod_{V\in\operatorname{Clop}(U)}\mathcal A(V).
\]
On nonempty connected $U$ this is $\mathcal A(U)_0$.
\end{proposition}
\begin{proof}
A topos is extensive, so products distribute over the indicated
coproducts. Define the operations on each summand: the old operations on
old inputs, the new point as additive identity, and as multiplicative
absorber. Every axiom follows on these summands. The two maps are
specified by sending the new point to old zero or recording its summand.
An inverse-image functor preserves finite limits and coproducts, so
preserves this construction. A section of the coproduct sheaf chooses a
locally constant summand, hence a clopen old-support subset $V$, and a
section on $V$; this description is reversible. A nonempty connected open
has just two such choices. On the empty open the sheaf has one section,
so that qualification is necessary.
\end{proof}

Every endomorphism of $A$ consequently extends while fixing the new zero.
This includes Frobenius actions on an arithmetic-site semiring sheaf.
Stalkwise, localization at its old zero gives $\B$ and at its new zero the
terminal semiring: the first localization identifies every old element
with $1$, while the second identifies $0$ and $1$. These are structural
statements about the given sheaf. They do not identify its stalks with a
fixed scalar set or supply an unstated morphism of arithmetic topoi.
